\section{Endpoint-distance coordinates for the same positive trace}
\label{sec:v103-gap}
The relative coordinate in \eqref{eq:v102-barycentric} need not be
preserved by a comparison. On a regular pairing chamber let
$e_-=a$, $e_+=b$, and put
\[
 x_-=u-a(\theta),\qquad x_+=b(\theta)-u,
 \qquad 0<x_\pm<\ell(\theta).
\]
These are distances in the roof variable, not spatial distances on
the billiard table. In either coordinate the original paired measure is
\begin{equation}\label{eq:v103-gap-measure}
 d\pi^\varepsilon=c^\varepsilon(\theta)\,d\theta\,dx_\pm.
\end{equation}
Indeed $x_-=s\ell$ and $x_+=(1-s)\ell$, so the factor $\ell$
in \eqref{eq:v102-pair-measure} is exactly the coordinate Jacobian.
No trace minimum, original weight or source selector is changed.

\begin{theorem}[Two endpoint-distance flows and their covariance]
\label{thm:v103-gap-coarea}
Fix $\Delta>0$. On the regular set
\[
 c^\varepsilon>0,\quad 0<x_\pm<\ell<\Delta,\quad e_\pm'\ne0,
\]
define $V_\pm=(e_\pm')^{-1}\partial_\theta$ at fixed $x_\pm$.
The corresponding stopped flows $\Theta_h^\pm$ have unit roof speed.
In $(\theta,s)$ coordinates their fields are
\begin{align}
 V_-&=\frac1{a'}\left(\partial_\theta-
                    \frac{s\ell'}\ell\partial_s\right),\notag\\
 V_+&=\frac1{b'}\left(\partial_\theta+
                    \frac{(1-s)\ell'}\ell\partial_s\right).
                                                    \label{eq:v103-gap-fields}
\end{align}
Their densities in $(t,x_\pm)$ and pullback Jacobians are
\begin{equation}\label{eq:v103-gap-jacobian}
 \rho_\pm^\varepsilon=\frac{c^\varepsilon}{|e_\pm'|},
 \qquad J_h^\pm(y)=
       \frac{\rho_\pm^\varepsilon(\Theta_h^\pm y)}
            {\rho_\pm^\varepsilon(y)}.
\end{equation}
Writing $S_h^\pm(t)$ for the stopped survivor fiber, the convention is
\begin{equation}\label{eq:v103-gap-covariance}
 \int_{S_h^\pm(t)}\varphi(\Theta_h^\pm y)J_h^\pm(y)
                          \,d\pi_t^\varepsilon(y)
 =\int_{\Theta_h^\pm S_h^\pm(t)}\varphi(z)
                          \,d\pi_{t+h}^\varepsilon(z).
\end{equation}
For fixed collision count, radius and exact label this holds jointly
for almost every $(t,h)$ and all nonnegative Borel integrands.
Each flow counts physical interiors with multiplicity one. The two
flows are distinct comparisons, not two copies of the source.
\end{theorem}
\begin{proof}
At fixed $x_-$ the roof is $a(\theta)+x_-$, and at fixed $x_+$
it is $b(\theta)-x_+$. In both cases its angular derivative is
$e_\pm'$. Thus $V_\pm\mathscr F=1$. Differentiating the two
identities for $x_\pm$ proves \eqref{eq:v103-gap-fields}.
The triangular coordinate change from $(\theta,x_\pm)$ to
$(t,x_\pm)$ has absolute determinant $|e_\pm'|$. Applying it to
\eqref{eq:v103-gap-measure} gives \eqref{eq:v103-gap-jacobian}.
In these coordinates the flow translates $t$ and fixes $x_\pm$;
its pullback density is therefore the displayed target-to-source
ratio, not its reciprocal.

On the open comparison domain the field is $C^1$. Its partial flow
and inverse are unique until a true comparison boundary is reached.
Apply change of variables on this one-to-one graph, then disintegrate
in $(t,h)$. The original countable physical charts and a countable
determining algebra identify one joint kernel version. Monotone
approximation proves \eqref{eq:v103-gap-covariance} for all
nonnegative integrands, also with Borel graph restrictions. Tonelli
sums the disjoint physical pieces. Smoothly compatible chart seams
are crossed by coordinate change; actual loss of pairing regularity
remains a stopping event. Neither $e_\pm'=0$ nor nonregular positive
source is discarded: its corresponding clock will be zero.
\end{proof}

\begin{proposition}[The nonvanishing trace at a regular critical entrance]
\label{prop:v103-critical-trace}
Suppose a chamber has a backward critical entrance with a regular
nongrazing stationary extension at $q_*$. Put $K_*=D_q^2F(q_*)$,
$e(\theta)=(\cos\theta,\sin\theta)$, and let $e_\perp$ be its
positive quarter-turn. Its incoming roof is $a=F(q_*)$, independent
of $\theta$, and its original weighted incoming trace is
\begin{equation}\label{eq:v103-critical-trace}
 f^\varepsilon(a+,\theta)=
 \frac{\omega(q_*)W^\varepsilon(q_*)}
      {4\pi c_*\,e_\perp(\theta)^{\mathsf T}K_*e_\perp(\theta)}.
\end{equation}
If its numerator is positive, then
\[
 |\partial_\theta\log f^\varepsilon(a+,\theta)|
                    \le 2\Lambda_\circ/\lambda_\circ.
\]
Nevertheless $c^\varepsilon$ is still the minimum with the outgoing
trace. This bound alone gives no relative bound for that minimum.
\end{proposition}
\begin{proof}
The positive Hessian gives a local inverse $q(g)$ near $g=0$.
All rays in this critical chamber approach the same $q_*$, so their
incoming roofs agree. At $g=ze$ direct differentiation gives
\[
 \det D_q(F,\arg g)=e_\perp^{\mathsf T}K(q)e_\perp.
\]
The optical factor and original guard have continuous values at the
regular nongrazing extension. Taking $z\downarrow0$ proves
\eqref{eq:v103-critical-trace}; no physical continuation beyond a
first-hit or section boundary is asserted. The derivative of the
quadratic form has magnitude at most $2\Lambda_\circ$, while the
form is at least $\lambda_\circ$. Division proves the bound.
A coincident singular extension not covered by the hypothesis remains
in the original source, not in the scope of this trace calculation.
\end{proof}

\begin{lemma}[Which zero-speed set can have paired mass]
\label{lem:v103-flat}
Within the regular chambers, the set $v=0$ with
$|a'|+|b'|>0$ is $\pi^\varepsilon$-null. Apart from transverse-null
sets, the remaining zero-speed source consists of angular intervals
on which both $a$ and $b$ are constant. On such an interval $I$,
its unselected roof density is exactly
\[
 \left(\int_Ic^\varepsilon(\theta)\,d\theta\right)
                  \one_{(a_I,b_I)}(t).
\]
No smallness of the sum of these flat-interval pulses is asserted.
\end{lemma}
\begin{proof}
For a fixed $\theta$ with $(a',b')\ne(0,0)$, the affine function
$(1-s)a'+sb'$ has at most one zero. The measure in $s$ is absolutely
continuous, proving the first assertion by Tonelli. The endpoint
maps are definable by Lemma~\ref{lem:v102-chambers}; their common
derivative-zero set is a finite union of points and intervals at
fixed count. The maps are constant on each interval. Integration
in the original $u$ coordinate gives the displayed density. Points
have zero transverse mass. The conclusion concerns exact zero
speed, not an estimate for neighborhoods of folds or stationary roofs.
\end{proof}
